Los resultados obtenidos verifican que el enfriamiento del sistema sigue de manera satisfactoria la ley de Newton, ya que la temperatura disminuye más rápido al inicio y luego lo hace cada vez más lentamente hasta aproximarse al equilibrio térmico con el ambiente. La linealidad de la gráfica logarítmica y el valor elevado de \(r^2\) muestran que el modelo describe bien el fenómeno observado. Las pequeñas desviaciones respecto al comportamiento ideal pueden atribuirse a incertidumbres experimentales como la precisión del termómetro, variaciones en la temperatura ambiente, errores de tiempo y condiciones no completamente uniformes de transferencia de calor.
